% Generated draft: requires mathematical and source review.
\documentclass[11pt]{article}
\usepackage[margin=25mm]{geometry}
\usepackage{fontspec}
\setmainfont{Latin Modern Roman}
\usepackage{amsmath,amssymb,mathtools,unicode-math}
\setmathfont{Latin Modern Math}
\usepackage{enumitem,needspace,xurl,hyperref}
\setlength{\parindent}{0pt}
\setlength{\parskip}{5pt}
\setlength{\emergencystretch}{3em}
\newcommand{\N}{\mathbb{N}}
\newcommand{\Z}{\mathbb{Z}}
\newcommand{\Q}{\mathbb{Q}}
\newcommand{\R}{\mathbb{R}}
\newcommand{\C}{\mathbb{C}}
\newcommand{\F}{\mathbb{F}}
\newcommand{\A}{\mathbb{A}}
\newcommand{\PP}{\mathbb{P}}
\newcommand{\E}{\mathbb{E}}
\newcommand{\eps}{\varepsilon}
\newcommand{\dd}{\,\mathrm{d}}
\begin{document}
\section*{UnsolvedMath ID: 30000938}
\textbf{Source problem code:} catalogue entry
\subsection{Definitions and mathematical statement}
Let \(\Lambda\subseteq\R^2\) be a fixed full-rank planar lattice and \(S\subseteq\R^2\) a fixed closed bounded convex region containing the origin. Under the usual homothetic-oval convention, put
\[
 S(r,P)=P+rS,\qquad r>0,\quad P\in\R^2,
\]
\[
 J(r,P)=\Lambda\cap S(r,P),\qquad
 N(r,P)=|J(r,P)|.
\]
The record refers to earlier ``same lattice and \(S\)-oval notation'' without explicitly defining \(S(r,P)\), so this homothety/translation interpretation requires confirmation. Closedness means that points on the boundary are included. No rotation of \(S\) is introduced.

Two configurations, viewed as finite subsets of \(\Lambda\), are equivalent precisely when
\[
 J\sim J'\quad\Longleftrightarrow\quad
 J'=J+v\text{ for some }v\in\Lambda.
\]
Only distinct subsets modulo lattice translation are counted, not parameter pairs \((r,P)\), which can yield the same subset. For a nonnegative integer \(n\), define
\[
 \mathcal C_n(\Lambda,S)=
 \{J(r,P):r>0,\ P\in\R^2,\ N(r,P)=n\}/\!\sim,
 \qquad K(n)=|\mathcal C_n(\Lambda,S)|.
\]
For a nonnegative integer \(N\), put
\[
 M(N)=\left|\{J(r,P):r>0,\ P\in\R^2,\ N(r,P)\leq N\}/\!\sim\right|
      =\sum_{n=0}^N K(n),
\]
when the empty configuration is included; a positive-weight convention omits \(n=0\). The questions ask to determine these counts as functions of the weight and cutoff for the fixed lattice and shape, including their intended growth/asymptotic behavior. No specific proposed formula, exponent, leading constant, or error term is supplied.

The original says translation by an integer vector, while the clean statement says a lattice vector. These agree after choosing a lattice basis identifying \(\Lambda\) with \(\Z^2\); they need not agree for literal ambient integer vectors without that identification.

The supplied assessment dated 2026-08-21 reports sharp bounds by Huxley--Zunic for quantities called \(L(N)\) and \(M(N)\) under a Quadrangle Condition, while general and circular cases retain unresolved asymptotic questions. Neither that condition nor \(L(N)\) is defined here, so \(L\) is not identified with the original \(K\), and the conditional bounds are not transferred to arbitrary \(S\).
\subsection{Short English statement}
Count the different lattice-point subsets cut out by translated, scaled copies of one fixed convex shape, identifying subsets that differ by a lattice translation. The problem asks both for a fixed number of points and for at most a given number, rather than counting the continuously varying placement parameters.
\subsection{Sources}
\begin{itemize}
\item Local UnsolvedMath registry, permanent ID \(30000938\), exact source code \texttt{catalogue entry}, proposed year 2008; no individual proposer is supplied.
\item \emph{Analytic Number Theory (2008)}, Oberwolfach Reports, DOI 10.4171/owr/2008/14, pp. 669--746. \url{https://doi.org/10.4171/owr/2008/14}.
\item Supplied reference: Huxley and Zunic, ``configurations II,'' \url{https://doi.org/10.1112/plms/pdt011}.
\item The supplied \texttt{corrected\_verified} note says the full statement was retained with notation placed in delimiters. The partially solved label and assessment checked 2026-08-21 are dated supplied triage requiring independent verification of the Quadrangle Condition, definitions, bounds, and remaining asymptotic scope.
\end{itemize}
\end{document}
